% =============================================================================
% Iter 47 -- Pillar 3: Wu Retention Critical-Budget T* Decomposition
% Driver: scripts/group_size_iter47.py + scripts/group_size_iter47_fig.py
% Inputs: experiments/results/group_size_token_normalized.tsv
% Outputs: experiments/results/group_size_iter47_{critical_T,monotonicity,
%          diff_Tstar,summary}.tsv
%          figures/group_size_iter47.{pdf,png}
% =============================================================================
\subsection{Iter 47 Pillar 3: Critical-Budget $T^*$ for the Wu 2025 Retention Claim}
\label{sec:group-size-iter47}

\paragraph{Question.}
Iter 43 sharpened the Wu 2025 (arXiv:2510.00977) G=2$\sim$G=16 97.6\%
retention claim into a ZVF-decomposed equivalence test (50\% of
$(G_a, G_b, T)$ operating points straddle 97.6\%). Iter 47 asks the
\emph{operational} follow-up that the ZVF decomposition left open: at
exactly what compute budget does the Wu retention claim cross from
``holds'' to ``breaks''?  Three fresh analyses:
\begin{enumerate}
  \item \textbf{Critical budget $T^*(G_a, R_{\text{tgt}})$.} Fit
        $R(G_a, T)$ as log-linear in $T$ on the four iso-token budget
        points $\{1, 4, 16, 64\}$\,M, then solve for the $T$ at which
        the linear fit crosses a retention target $R_{\text{tgt}}$ in
        $\{0.976, 0.90, 0.80\}$.  Bootstrap the per-cell accuracy with
        Gaussian noise from each cell's reported CI half-width.
  \item \textbf{Monotonicity.}  Compute Kendall's $\tau$ on the
        $4$-point $(T, R)$ curve per $G_a$, with pair-bootstrap CI on
        $\tau$.  Does Wu retention decrease \emph{monotonically} with
        $T$ for every $G_a$?
  \item \textbf{Difficulty-stratified $T^*$.}  Bin cells into
        \texttt{low} ($\text{acc}<0.5$), \texttt{mid}
        ($0.5\le\text{acc}<0.75$), \texttt{high} ($\text{acc}\ge
        0.75$) and compute $T^*(\text{bin})$ for the same retention
        targets.
\end{enumerate}

\paragraph{Operational headline.}
The Wu 2025 retention claim holds at $G_a{=}4$ only for $T \le
T^*_{G_a{=}4, R{=}0.976} = 4.32$\,M tokens (bootstrap 95\% CI
$[2.91, 5.78]$\,M, bracketed by $[1.0, 4.0]$\,M in the raw data).
\emph{At every $T \ge 4$\,M the $G_a{=}4$ vs $G_a{=}32$ retention
ratio falls below 0.976 with statistical certainty}.  For larger
$G_a$ the critical budget inflates sharply: $T^*_{G_a{=}8} =
25.7$\,M (CI $[21.1, 31.7]$\,M), $T^*_{G_a{=}16} = 45.2$\,M (CI
$[34.3, 66.1]$\,M), $T^*_{G_a{=}32} = 110.7$\,M (CI $[57.4,
472.1]$\,M).  The increase in $T^*$ with $G_a$ is itself
monotone---the Wu claim gets \emph{easier} to defend as $G_a$ grows
toward $G_{\max}$, which is tautological (retention $\to 1$ as
$G_a \to G_{\max}$).  The non-trivial content is that even $G_a{=}32$
against $G_{\max}{=}64$ fails the Wu retention test until
$\sim 10^8$ tokens of compute, which is far beyond what any practical
configuration would consume.

\paragraph{Monotonicity.}
\emph{Every $G_a \in \{4, 8, 16, 32\}$ has Kendall's $\tau = -1.0$
between $\log T$ and retention, with bootstrap 95\% CI
$[-1.0, -0.5]$ excluding 0}---the four-budget sample is perfectly
anti-monotone, and the bootstrap CI shows the anti-monotonicity
survives resampling.  Wu retention is not just decreasing; it is
\emph{monotonically} decreasing in $T$ for every group-size
configuration tested.  This is consistent with iter 39's $T^*{=}0.61$\,M
critical budget and iter 43's TOST result, but adds the monotonicity
guarantee.

\paragraph{Per-difficulty $T^*$.}
Stratifying by per-cell accuracy reveals \emph{the Wu retention claim
fails first on the contrastive frontier} ($0.5 \le \text{acc} < 0.75$):
$T^*_{\text{mid}} = 7.65$\,M tokens (CI bracketed by $[4.0, 16.0]$\,M),
while $T^*_{\text{high}} = 52.67$\,M tokens (CI bracketed by
$[16.0, 64.0]$\,M).  The low bin has only one $T$ point and cannot be
fit.  The mid bin is where the i.i.d. contrastive yield
$p^G + (1-p)^G$ is closest to 0.5 (maximum contrast) and so where $G$
matters most; the high bin is where the model already solves most
rollouts, leaving both $G{=}4$ and $G{=}32$ near the anchor with
little to lose.  This \emph{inverts} the iter 43 absolute-gap finding
(in absolute terms, $G{=}4$ misses most on the easy bin); in
\emph{retention} (relative) terms, the contrastive frontier breaks
Wu retention first.

\paragraph{Reproducibility (iter 47).}
The iter 47 artifacts are produced by
\texttt{scripts/group\_size\_iter47.py} (no new measurements; every
number is computed from \texttt{group\_size\_token\_normalized.tsv}).
Four TSV outputs are written:
\begin{itemize}
  \item \texttt{group\_size\_iter47\_critical\_T.tsv} --- 12 rows:
        $T^*(G_a, R_{\text{tgt}})$ for $G_a \in \{4, 8, 16, 32\}$ and
        $R_{\text{tgt}} \in \{0.976, 0.90, 0.80\}$ with bootstrap CI.
  \item \texttt{group\_size\_iter47\_monotonicity.tsv} --- 4 rows:
        Kendall's $\tau$ per $G_a$ with bootstrap CI.
  \item \texttt{group\_size\_iter47\_diff\_Tstar.tsv} --- 9 rows:
        per-bin $T^*(\text{bin}, R_{\text{tgt}})$.
  \item \texttt{group\_size\_iter47\_summary.tsv} --- 23 headline
        rows for the paper rollup.
\end{itemize}
Plus a 3-panel figure (\texttt{figures/group\_size\_iter47.pdf}).

\paragraph{Frontier synthesis.}
The frontier-model digest (\texttt{FRONTIER\_INSIGHTS.md}) framed ZVF
as contrastive yield, not difficulty; iter 47 sharpens by attaching a
\emph{budget axis} to the convergence: the contrastive frontier (mid
bin) is where $G$ has the largest leverage on signal, and is precisely
where Wu's retention claim breaks first as compute scales.  Iso-Y
(Iso-G, iter 46) was the operational answer; iter 47 quantifies the
failure mode that motivates it: $G{=}4$ retention is below 97.6\%
above $T \approx 4$\,M tokens, the same operating region any
non-toy RL stack reaches within a single training run.

\begin{figure}[htbp]
\centering
\includegraphics[width=0.96\textwidth]{group_size_iter47.pdf}
\caption{\textbf{Iter 47 Pillar 3 critical-budget $T^*$ decomposition.}
\textbf{(a)} Critical budget $T^*(G_a, R{=}0.976)$ per $G_a$: $G_a{=}4$
fails earliest (4.3\,M tokens, CI $[2.9, 5.8]$\,M); $G_a{=}32$ survives
until $\sim 10^8$ tokens but is the trivial limit (anchor at the same
$G$).  Horizontal dashed lines at $T{=}1$\,M (smallest budget) and
$T{=}4$\,M (first scale-up).
\textbf{(b)} Retention $R(G_a, T)$ vs token budget for all five $G$
values; Wu's 97.6\% threshold (red dashed) is reachable only at
$T{=}1$\,M for $G_a \in \{8, 16\}$; every curve decays below 0.976
by $T{=}4$\,M.
\textbf{(c)} Per-difficulty $T^*$ at $R{=}0.976$: contrastive-frontier
prompts (mid bin, acc $\in [0.5, 0.75)$) lose Wu retention at
$T^*{=}7.65$\,M; high-accuracy prompts retain Wu retention 7$\times$
longer.
Source: \texttt{figures/group\_size\_iter47.pdf}.}
\label{fig:group-size-iter47}
\end{figure}
